\section{Proof of Correctness on White Noise}
In this section, we prove that the expected DDL of white noise under a GMM is a monotonically increasing function of the number of components in the GMM. 

\begin{lemma} \label{lemma:r}
When clustering white noise on $[0,1]^m$, with a $k$-component GMM, the radius $r$ of each cluster is approximated by 
 \begin{equation} \label{eq:r}
\frac{1}{\sqrt{3}\sqrt[m]{k}}\left(\frac{2}{\sqrt{3}}\right)^{1/m} \,.
\end{equation}
\end{lemma}
\begin{proof}
For balls of radius $r$, the distance along each coordinate axis between their centres will be $2r$ for one of the dimensions and $2r\frac{\sqrt{3}}{2}$ for all other dimensions. This is because the centres of each 3 touching balls will form the vertices of an equilateral triangle with side length $2r$, which will then have height $2r\frac{\sqrt{3}}{2}$. Thus, the number of balls of radius $r$ that can fit inside each axis is $\frac{1}{2r}$ for the first axis and $\frac{1}{2r}\frac{2}{\sqrt{3}}$ for all other axes. The total number of balls that can fit inside $[0,1]^m$ is, therefore 
\[
\frac{1}{(2r)^m}\left(\frac{2}{\sqrt{3}}\right)^{m-1}\,.
\]
Conversely, given that the GMM will have $k$ clusters, we can approximate the radius of each cluster using 
\begin{gather*}
\frac{1}{(2r)^m}\left(\frac{2}{\sqrt{3}}\right)^{m-1} = k \\
(2r)^m = \frac{1}{k}\left(\frac{2}{\sqrt{3}}\right)^{m-1}  \\
2r = \frac{1}{\sqrt[m]{k}}\left(\frac{2}{\sqrt{3}}\right)^{\frac{m-1}{m}}  \\
r = \frac{1}{\sqrt{3}\sqrt[m]{k}}\left(\frac{\sqrt{3}}{2}\right)^{1/m}  \\
\end{gather*}
\end{proof}

\begin{proposition} \label{proposition:m-ball-pdf}
For uniformly distributed points in an $m$-dimensional hyperball of radius $r$, the pdf of the distance of a point from the centre is given by
\[
p(x) = m \frac{x^{m-1}}{r^m}
\]
\end{proposition}
\begin{proof}
The probability density is clearly proportional to $x^{m-1}$, thus we need only to find the constant $c$ such that the pdf is appropriately normalized. 
\begin{gather*}
    c\int_{0}^r x^{m-1}dx = 1 \\
    c \frac{x^m}{m} \big |_0^r = 1\\
    c \frac{r^m}{m} = 1\\
    c = \frac{m}{r^m} \,.
\end{gather*}
\end{proof}

\begin{lemma} \label{lemma:a}
 When clustering white noise in $[0,1]^m$ dimensions, with a $k$-component GMM, the expected squared distance of a point from the centroid of its cluster, denoted $a$, is given by
 \begin{equation} \label{eq:a}
a = \frac{m}{3(m+2)}\left(\frac{\sqrt{3}}{2}\right)^{2/m}\frac{1}{k^{2/m}}\,,
\end{equation}
\end{lemma}
\begin{proof}
Using Proposition \ref{proposition:m-ball-pdf}, the expected squared distance from the centroid can be calculated directly from the pdf as
\begin{gather*}
    a = \frac{m}{r^m}\int_{0}^r x^{m+1}dx \\
    a = \frac{m}{r^m} \frac{x^{m+2}}{m+2} \big |_0^r = 1\\
    a = r^2\frac{m}{m+2}
\end{gather*}
Substituting $r$ from Lemma \ref{lemma:r} we get
\[
a = \frac{m}{3(m+2)}\left(\frac{\sqrt{3}}{2}\right)^{2/m}\frac{1}{k^{2/m}}\,,
\]
as desired.
\end{proof}

\begin{lemma} \label{lemma:variance}
When clustering white noise in $[0,1]^m$ dimensions, with a $k$-component GMM, the fit covariance matrix will $\Sigma$ will satisfy $\Sigma = \sigma I$, where 
\[
\sigma = \frac{1}{3(m+2)}\left(\frac{\sqrt{3}}{2}\right)^{2/m}\frac{1}{k^{2/m}}
\]
\end{lemma}

\begin{proof}
It is clear that $\Sigma$ will be of the form $\sigma I$ for some $\sigma$, as the clusters will all be identical by symmetry (up to approximation at the boundary of the hyperbox, but this is small for more than a few clusters). 

Next, observe that the expected squared distance of a point from the centroid of its cluster is, by linearity of expectation, equal to the sum of the expected squared distances in each coordinate, i.e. $m\sigma$. Then, by Lemma \ref{lemma:a}, we get 
\[
\sigma = \frac{1}{3(m+2)}\left(\frac{\sqrt{3}}{2}\right)^{2/m}\frac{1}{k^{2/m}}
\]
as desired.
\end{proof}

\begin{proposition} \label{proposition:fx}
The DDL of a point under the distribution of its cluster depends only on its distance from the centroid of its cluster.
\end{proposition}
\begin{proof}

The probability density of a point $z$ under a cluster with centroid $\mu$ is
\begin{align*}
p(z)    =& \frac{1}{\sqrt{2\pi |\Sigma|}}\exp{\left(\frac{-1}{2}(z-\mu)^T\Sigma^{-1}(z-\mu)\right)} \\
        =& \frac{1}{\sqrt{2\pi |\Sigma|}}\exp{\left(\frac{-|(z-\mu)|^2}{2\sigma}\right)} \,.
\end{align*}
Thus, the DDL under the cluster distribution, which we denote $\bar{D}$, is given by
\begin{gather*}
 -\ln{\left(\frac{1}{\sqrt{2\pi |\Sigma|}}\exp{\left(\frac{-|(z-\mu)|^2}{2\sigma}\right)}\right)} \\
        = \frac{1}{2} \ln{(2\pi |\Sigma|}) + \frac{|(z-\mu)|^2}{2\sigma} \,.
\end{gather*}
In what follows, we will use the function $f(x)$ to denote the DDL under the cluster distribution of a point a distance $x$ from its centroid, where
\[
f(x) = \frac{1}{2} \ln{(2\pi |\Sigma|}) + \frac{x^2}{2\sigma} \,.
\]
\end{proof}

\begin{definition}
Denote as outliers, those points with greater DDL under their cluster distribution than under the prior distribution.
\end{definition}
 
%To compute $\bar{D}$, we consider separately the cases where there are outliers and where there are no outliers. 
\begin{lemma} \label{lemma:d}
When clustering white noise, so that the prior distribution is $[0,1]^m$, a point is an outlier if and only if it is greater than a distance $d$ from is centroid, where $d=\sqrt{\sigma \ln{\frac{1}{2\pi |\Sigma|}}}$.
\end{lemma}
\begin{proof}
The DDL of a point treated as an outlier on the uniform $m$-box $[0,1]^m$ is 0, because treating as an outlier means using the prior distribution, which is uniform $p(x)=1$, giving DDL of $\ln{1} = 0$. Thus, a point a distance $x$ from its centroid is not an outlier if and only if its DDL under the distribution of its cluster is strictly negative:
\begin{gather*}
    DDL(x) < 0 \\
    \frac{1}{2} \ln{(2\pi |\Sigma|}) + \frac{x^2}{2\sigma} < 0 \\
    \frac{x^2}{\sigma} < -\ln{(2\pi |\Sigma|}) \\
    x^2 < \sigma \ln{\frac{1}{2\pi |\Sigma|}} \\
    x < \sqrt{\sigma \ln{\frac{1}{2\pi |\Sigma|}}} \,.
\end{gather*}
We refer to this distance $d$ as the inlier radius.
\end{proof}

\begin{lemma} \label{lemma:saturating-k}
When clustering white noise in $m$ dimensions using a GMM with $k$ components, some points will be classed as outliers if and only if $k$ satisfies
\[
k < \frac{2e\sqrt{\pi}}{\sqrt{3}}\left(\frac{e}{3(m+2)}\right)^{m/2} \\
\]
\end{lemma} \label{lemma:k-at-which-contained}
\begin{proof}
A point will be an outlier if and only if it is within the cluster, i.e. with a distance $r$ from its centroid, but outsider the inlier radius $d$. This is possible if and only if 
\begin{gather*}
    d < r \\
    \sqrt{\sigma \ln{\frac{1}{2\pi |\Sigma|}}} <\frac{1}{\sqrt{3}\sqrt[m]{k}}\left(\frac{\sqrt{3}}{2}\right)^{1/m} \\
    \sigma \ln{\frac{1}{2\pi |\Sigma|}} < \frac{1}{3k^{2/m}}\left(\frac{\sqrt{3}}{2}\right)^{2/m} \,.
\end{gather*}
As $\Sigma = \sigma I$, we can sub in $|\Sigma| = \sigma^m$, and also sub in $\sigma$ from Lemma \ref{lemma:variance}:
\begin{gather*}
    \frac{1}{3(m+2)}\left(\frac{\sqrt{3}}{2}\right)^{2/m}\frac{1}{k^{2/m}}\ln{\frac{1}{2\pi \sigma^m}} < \frac{1}{3k^{2/m}}\left(\frac{\sqrt{3}}{2}\right)^{2/m} \\
    \frac{1}{(m+2)}\ln{\frac{1}{2\pi \sigma^m}} < 1 \\
    \ln{\frac{1}{2\pi \sigma^m}} < m+2 \\
    \ln{\frac{1}{\sigma^m}} < m+2 + \ln 2{(\pi)} \\
    m\ln{\frac{1}{\sigma}} < m+2 + \ln 2{(\pi)} \\
    \ln{\frac{1}{\sigma}} < 1 + \frac{2 + \ln 2{(\pi)}}{m}   \\
    \sigma > \frac{1}{e^{1 + \frac{2 + \ln 2{(\pi)}}{m}}}   \\
    \frac{1}{3(m+2)}\left(\frac{\sqrt{3}}{2}\right)^{2/m}\frac{1}{k^{\frac{2m-2}{m}}} > \frac{1}{\sqrt[m]{\pi}e^{1 + \frac{2}{m}}}   \\
    \frac{\sqrt[m]{\pi}e^{1 + \frac{2}{m}}}{3(m+2)}\left(\frac{\sqrt{3}}{2}\right)^{2/m} > k^{2/m} \\
    k^{2/m} < \frac{\sqrt[m]{\pi}e^{1 + \frac{2}{m}}}{3(m+2)}\left(\frac{\sqrt{3}}{2}\right)^{2/m} \\
    k < \frac{\sqrt{\pi}e^{\frac{m}{2} + 1}}{(3(m+2))^{m/2}}\left(\frac{2}{\sqrt{3}}\right) \\
    k < \frac{\sqrt{\pi}e^{\frac{m}{2} + 1}}{(3(m+2))^{m/2}}\frac{\sqrt{3}}{2} \\
    k < \frac{\sqrt{\pi}e\sqrt{3}}{2}\frac{e^{m/2}}{(3(m+2))^{m/2}} \\
    k < \frac{e\sqrt{3\pi}}{2}\left(\frac{e}{3(m+2)}\right)^{m/2} \\
\end{gather*}
\end{proof}

\begin{lemma} \label{lemma:form-of-EDDL}
For a $k$-component GMM fit on $m$-dimensional white noise, the expected DDL of a point is 
\[
\frac{m}{r^m}\int_0^x x^{m-1} \min(\{0,f(x)\}) dx + \ln{k} \,.
\]
with $f(x)$ defined as in Proposition \ref{proposition:fx}.
\end{lemma}
\begin{proof}
The DDL of a point can be decomposed as the number of bits needed to specify which cluster the point belongs to, plus the DDL of the point under that cluster. The former is equal to $\ln{k}$, because each cluster will, by symmetry, be equally sized. The latter, we denote $\bar{D}$, and is given by
\[
\bar{D} = \int_0^r p(r) \min(\{0,f(x)\}) \,,
\]
where $p(r)$ is the probability density function of the distance of a point from its centroid and $f(x)$ specifies the DDL of a point in terms of the distance from its centroid. 

Substituting for $p(x)$ from Proposition \ref{proposition:m-ball-pdf} gives 
\[
\bar{D} = \frac{m}{r^m}\int_0^x x^{m-1} \min(\{0,f(x)\}) dx \,.
\]
%As we define outliers to be those points with greater DDL under their cluster distribution than under the prior distribution, and as the DDL under the prior distribution is $0$, we have that $g(x) = \min(\{0,f(x)\})$, with $f(x)$ defined as in Proposition \ref{proposition:fx}.
\end{proof}

\begin{lemma} \label{lemma:increasing-when-outliers}
When clustering white noise on $[0,1]^m$ with a $k$-component GMM, for $k \leq d$, the expected DDL is an increasing function of $k$. 
\end{lemma}
\begin{proof}
When some points are outliers, we have, using Lemma \ref{lemma:form-of-EDDL}
\begin{align*}
    \bar{D} &= \frac{m}{r^m}\int_0^r x^{m-1} \min(\{0,f(x)\}) dx \\
            &= \frac{m}{r^m}\int_{\{x \in [0,r] | f(x) < 0\}} r^{m-1} f(x) dx \,.
\end{align*}
By Lemma \ref{lemma:d}, $\{x \in [0,r] | f(x) < 0\} = [0,d)$, giving
\begin{align*}
     \bar{D} &= \frac{m}{r^m}\int_0^d x^{m-1} f(x) dx \\
\end{align*}
Substituting $f(x)$ from Proposition \ref{proposition:fx} and integrating:
\begin{align*}
    \bar{D} &= \frac{m}{r^m}\int_0^d x^{m-1} \left(\frac{1}{2} \ln{(2\pi |\Sigma|}) + \frac{x^2}{2\sigma}\right)  dx \\
            &= \frac{m}{r^m}\int_0^d \frac{1}{2} \ln{(2\pi |\Sigma|})x^{m-1}  + \frac{x^{m+1}}{2\sigma}  dx \\
            &= \frac{m}{r^m} \left(d^{m} \frac{1}{2m} \ln{(2\pi |\Sigma|}) + \frac{d^{m+2}}{2\sigma(m+2)} \right) \\ 
            &= \frac{d^{m}}{2r^m}  \left(\ln{(2\pi |\Sigma|}) + d^{2}\frac{m}{\sigma(m+2)} \right) \,.
\end{align*}
Substituting $d$ from Lemma \ref{lemma:d}:
\begin{align*}
     \bar{D} &= \frac{(\sigma \ln{\frac{1}{2\pi |\Sigma|}})^{m/2}}{2r^m}  \left(\ln{(2\pi |\Sigma|}) + (\sigma \ln{\frac{1}{2\pi |\Sigma|}})\frac{m}{\sigma(m+2)} \right) \\
            &= \frac{(\sigma \ln{\frac{1}{2\pi |\Sigma|}})^{m/2}}{2r^m}  \left(\ln{(2\pi |\Sigma|}) + (\ln{\frac{1}{2\pi |\Sigma|}})\frac{m}{m+2} \right) \\
            &= \frac{(\sigma \ln{\frac{1}{2\pi |\Sigma|}})^{m/2}}{2r^m}  \left(\ln{\frac{1}{2\pi |\Sigma|}}\frac{-2}{m+2} \right) \\
            &= \frac{-\sigma ^{m/2}}{r^m(m+2)}  \left(\ln{\frac{1}{2\pi |\Sigma|}} \right)^{\frac{m}{2}+1} \,.
\end{align*}
Substituting $\sigma$ from Lemma \ref{lemma:variance}:
\begin{align*}
      \bar{D} &= -\left(\frac{1}{3(m+2)}\left(\frac{\sqrt{3}}{2}\right)^{2/m}\frac{1}{k^{2/m}}\right)^{m/2} \frac1{r^m(m+2)} \left(\ln{\frac{1}{2\pi |\Sigma|}} \right)^{\frac{m}{2}+1} \\
            &= -\left(\frac{1}{(3(m+2))^{m/2}}\left(\frac{\sqrt{3}}{2}\right)\frac{1}{k}\right) \frac{1}{r^m(m+2)} \left(\ln{\frac{1}{2\pi |\Sigma|}} \right)^{\frac{m}{2}+1} \\
            &= -\left(\frac{1}{(3(m+2))^{m/2}}\right) \left(\frac{\sqrt{3}}{2}\right) \frac{1}{kr^m(m+2)} \left(\ln{\frac{1}{2\pi |\Sigma|}} \right)^{\frac{m}{2}+1} \\
            &=  \frac{-3 \sqrt{3}}{r^m2(3)^{m/2}k(m+2)} \left(\frac{1}{(m+2)}\ln{\frac{1}{2\pi |\Sigma|}} \right)^{\frac{m}{2}+1} \,.
\end{align*}
Substituting $r$ from Lemma \ref{lemma:r}:
\begin{align*}
       \bar{D} &= -\left(\sqrt{3}\sqrt[m]{k}\left(\frac{2}{\sqrt{3}}\right)^{1/m} \right)^m\frac{-3 \sqrt{3}}{2(3)^{m/2}k(m+2)} \left(\frac{1}{(m+2)}\ln{\frac{1}{2\pi |\Sigma|}} \right)^{\frac{m}{2}+1} \\
       \bar{D} &= -(\sqrt{3})^m k\frac{2}{\sqrt{3}} \frac{-3 \sqrt{3}}{2(3)^{m/2}k(m+2)} \left(\frac{1}{(m+2)}\ln{\frac{1}{2\pi |\Sigma|}} \right)^{\frac{m}{2}+1} \\
       \bar{D} &= \frac{-3}{(m+2)^{\frac{m}{2}+2}} \left(\ln{\frac{1}{2\pi |\Sigma|}} \right)^{\frac{m}{2}+1} \,.
\end{align*}
Substituting $\Sigma=\sigma I$, with $\sigma$ from Lemma \ref{lemma:variance}:
\begin{align*}
       \bar{D} &= \frac{-3}{(m+2)^{\frac{m}{2}+2}} \left(\ln{\frac{1}{2\pi}} + \ln{\frac{1}{|\Sigma|}} \right)^{\frac{m}{2}+1} \\
            &= \frac{-3}{(m+2)^{\frac{m}{2}+2}} \left(\ln{\frac{1}{2\pi}} + \ln{\frac{1}{\sigma^m}} \right)^{\frac{m}{2}+1} \\
            &= \frac{-3}{(m+2)^{\frac{m}{2}+2}} \left(\ln{\frac{1}{2\pi}} + \ln{\left(3^m(m+2)^m\frac{4}{3}k^2\right)} \right)^{\frac{m}{2}+1} \\
            &= \frac{-3}{(m+2)^{\frac{m}{2}+2}} \left((m)\ln{3} +m\ln{(m+2)} + \ln{\tfrac{4}{3}} + 2 \ln{k} - \ln{2\pi} \right)^{m/2} \\
            &= \frac{-3}{(m+2)^{\frac{m}{2}+2}} \left(m\ln{(3m+6)} + 2 \ln{k} + \ln{\tfrac{2}{3\pi}} \right)^{\frac{m}{2}+1} \\
            &= \frac{-3}{(m+2)^{\frac{m}{2}+2}} \left(m\ln{(3m+6)} + 2 \ln{k} + \ln{\tfrac{2}{3\pi}} \right)^{\frac{m}{2}+1} \\
\end{align*}
Differentiating, with respect to $k$, the total DDL of $\bar{D} + \ln{k}$:
\begin{align*}
    &\frac{d}{dk}\left(\frac{-3}{(m+2)^{\frac{m}{2}+2}}(m \ln{(3m+6)}+\ln{\tfrac{2}{3\pi}}+2\ln{k})^{\frac{m}{2}+1} + \ln{k}\right) \\
    &=\frac{-3}{(m+2)^{\frac{m}{2}+2}}\frac{d}{dk}\left((m \ln{(3m+6)}+\ln{\tfrac{2}{3\pi}}+2\ln{k})^{\frac{m}{2}+1}\right) + \frac{1}{k} \\
    &=\frac{-3}{(m+2)^{\frac{m}{2}+2}}\left((\frac{m}{2}+1)(m \ln{(3m+6)}+\ln{\tfrac{2}{3\pi}}+2\ln{k})^{m/2}\frac{2}{k}\right) + \frac{1}{k} \\
    &=\frac{-3}{(m+2)^{\frac{m}{2}+2}}\frac{(m+2)}{k}\left(m \ln{(3m+6)}+\ln{\tfrac{2}{3\pi}}+2\ln{k}\right)^{m/2} + \frac{1}{k} \\
    &=\frac{-3}{k(m+2)^{\frac{m}{2}+1}}\left(m \ln{(3m+6)}+\ln{\tfrac{2}{3\pi}}+2\ln{k}\right)^{m/2} + \frac{1}{k} \,.
\end{align*}
We want to show this derivative is positive:
\begin{align*}
    &\frac{-3}{k(m+2)^{\frac{m}{2}+1}}^{\frac{m}{2}+2}\left(m \ln{(3m+6)}+\ln{\tfrac{2}{3\pi}}+2\ln{k}\right)^{m/2} + \frac{1}{k} > 0\\
    &\iff \frac{1}{k} > \frac{3}{k(m+2)^{\frac{m}{2}+1}}\left(m \ln{(3m+6)}+\ln{\tfrac{2}{3\pi}}+2\ln{k}\right)^{m/2}\\
    &\iff \frac{3}{(m+2)^{\frac{m}{2}+1}}\left(m \ln{(3m+6)}+\ln{\tfrac{2}{3\pi}}+2\ln{k}\right)^{m/2} < 1\\
\end{align*}
As, by assumption, some points are outliers, we can use the upper bound on $k$ from Lemma \ref{lemma:k-at-which-contained}:
\begin{align*}
    &\frac{3}{(m+2)^{\frac{m}{2}+1}}\left(m \ln{(3m+6)}+\ln{\tfrac{2}{3\pi}}+2\ln{k}\right)^{m/2} \\
    &= \frac{3}{(m+2)^{\frac{m}{2}+1}}\left(m \ln{(3m+6)}+\ln{\tfrac{2}{3\pi}}+\ln{k^2}\right)^{m/2} \\
    &\leq \frac{3}{(m+2)^{\frac{m}{2}+1}}\left(m \ln{(3m+6)}+\ln{\tfrac{2}{3\pi}}+\ln{\left(\frac{4\pi e^2}{3}\left(\frac{e}{3(m+2)}\right)^m\right)}\right)^{m/2} \\
    &= \frac{3}{(m+2)^{\frac{m}{2}+1}}\left(m \ln{(3m+6)} + \ln{\tfrac{2}{3\pi}} + 2 + \ln{\tfrac{4\pi}{3}} + m - m\ln{(3m+6)}\right)^{m/2} \\
    &= \frac{3}{(m+2)^{\frac{m}{2}+1}}\left(2 + \ln{\tfrac{2}{3\pi}} + \ln{\tfrac{4\pi}{3}} + m\right)^{m/2} \\
    &= \frac{3}{(m+2)^{\frac{m}{2}+1}}\left(m + 2 + \ln{\tfrac{8}{9}}\right)^{m/2} \\
    &< \frac{3}{(m+2)^{\frac{m}{2}+1}}(m + 2)^{m/2} \\
    &= \frac{3}{m+2} \leq 1 \,,
\end{align*}
where the last inequality holds because the dimension $m$, is $\geq 1$.
\end{proof}

\begin{lemma} \label{lemma:constant-when-no-outliers}
When clustering white noise on $[0,1]^m$ with a $k$-component GMM, for $k>d$, the expected DDL is independent of $k$. 
\end{lemma}
\begin{proof}
When no points are outliers, $f(x) = g(x)$ for all $x$, so, using Lemma \ref{lemma:form-of-EDDL}
\begin{align*}
    \bar{D} &= \frac{m}{r^m}\int_0^r x^{m-1} f(x) dx \\
            &= \frac{m}{r^m}\int_0^r x^{m-1} \left(\frac{1}{2} \ln{(2\pi |\Sigma|}) + \frac{x^2}{2\sigma}\right)  dx \\
            &= \frac{m}{r^m} \left(\frac{r^m}{2m} \ln{(2\pi |\Sigma|}) + \frac{r^{m+2}}{2\sigma(m+2)}\right) \\
            &= \frac{1}{2} \ln{(2\pi |\Sigma|}) + \frac{r^2 m}{\sigma(m+2)} \,.
\end{align*}
Substituting $r$ from Lemma \ref{lemma:r}:
\begin{align*}
    \bar{D} &= \frac{1}{2} \left(\ln{(2\pi |\Sigma|}) + \frac{1}{3k^{2/m}}\left(\frac{\sqrt{3}}{2}\right)^{2/m}\frac{m}{\sigma(m+2)}\right)  dx \\
\end{align*}
Substituting $\sigma$ from Lemma \ref{lemma:variance}:
\begin{align*}
    \bar{D} &= \frac{1}{2} \left(\ln{(2\pi |\Sigma|}) + \frac{1}{3k^{2/m}}\left(\frac{\sqrt{3}}{2}\right)^{2/m}3(m+2)\left(\frac{2}{\sqrt{3}}\right)^{2/m}k^{2/m}\frac{m}{m+2}\right)  dx \\
    \bar{D} &= \frac{1}{2} \left(\ln{(2\pi |\Sigma|}) + m \right) \,.
\end{align*}
Substituting $\Sigma=\sigma I$, with $\sigma$ from Lemma \ref{lemma:variance}:
\begin{align*}
    \bar{D} &= \frac{1}{2} \left(\ln{(2\pi |\Sigma|}) + m \right) \\
            &= \frac{1}{2} \left(\ln{(2\pi)} + m \ln{{\sigma}} + m \right) \\
            &= \frac{1}{2} \left(\ln{(2\pi)} + m \ln{\left(\frac{1}{3(m+2)}\left(\frac{\sqrt{3}}{2}\right)^{2/m}\frac{1}{k^{2/m}}\right)} + m \right) \\
            &= \frac{1}{2} \left(\ln{(2\pi)} - m \ln{\left(3(m+2)\left(\frac{2}{\sqrt{3}}\right)^{2/m}k^{2/m}\right)} + m \right) \\
            &= \frac{1}{2} \left(\ln{(2\pi)} - m \ln{(3m+6)} + 2\ln{\tfrac{2}{\sqrt{3}}} -2\ln{k} + m \right) \\
            &= \frac{1}{2} \left(\ln{\frac{4\pi}{\sqrt{3}}}+ m \ln{\left(1+\frac{1}{3m+6}\right)} - 2\ln{k} + m \right) \\
\end{align*}
The full DDL is then
\begin{align*}
    & \bar{D} + \ln{k} \\
    &= \frac{1}{2} \left(\ln{\frac{4\pi}{\sqrt{3}}}+ m \ln{\left(1+\frac{1}{3m+6}\right)} - 2\ln{k} + m \right)  + \ln{k} \\
    &= \frac{1}{2} \left(\ln{\frac{4\pi}{\sqrt{3}}}+ m \ln{\left(1+\frac{1}{3m+6}\right)} + m \right) \\
\end{align*}
\end{proof}


\begin{theorem}
When clustering white noise in $[0,1]^m$, using a GMM with $k$ components, the expected DDL of a point is a monotonically increasing function of $k$.
\end{theorem}
\begin{proof}
By Lemma \ref{lemma:increasing-when-outliers}, the expected DDL is strictly increasing in $k$ up to $k=d$. By Lemma \ref{lemma:constant-when-no-outliers}, the expected DDL is constant in $k$ for $k > d$.
\end{proof}




\section{Worked Examples}
Section \ref{subsec:worked-example} presented a single worked example of my image complexity metric on a single, randomly chosen image from ImageNet. This section contains further examples from ImageNet and others from Cifar, MNIST, and DTD2.

\begin{figure*}
    \centering
    \begin{tabular}{@{}c@{}c@{}}
    \includegraphics[width=0.53\textwidth]{imgs/chainsaw_im.JPEG} & 
    \end{tabular}
    \caption{ \small  \small Example of a relatively high-resolution real-world image from imagenet. ID: n03445777\_10762.} \label{fig:imgs/im-chainsaw-image}
\end{figure*}
\noindent \begin{large} \textbf{Layer 1} \end{large} \\
Number of points to be clustered (pixels): 50246 \\
Number of components found by MDL, as per \eqref{eq:k-star}: 8 \\
Assign each pixel a label from $0, \dots, 6$, and form patch signatures as multisets of labels inside all $4 \times 4$ patches, which gives 48450 patches, of 2628 different unique values. \\
Entropy of resulting categorical distribution of patch signatures: \textbf{9.280} \\

\noindent \begin{large} \textbf{Layer 2} \end{large} \\
Number of points to be clustered: 48450 \\
Number of components found by MDL, as per \eqref{eq:k-star}: 6 \\
Assign each point a label from $0, \dots, 7$, and form patch signatures as multisets of labels inside all $8 \times 8$ patches, which gives 44954 patches, of 3809 different unique values. \\
Entropy of resulting categorical distribution of patch signatures: \textbf{9.363} \\

\noindent \begin{large} \textbf{Layer 3} \end{large} \\
Number of points to be clustered: 44954 \\
Number of components found by MDL, as per \eqref{eq:k-star}: 8 \\
Assign each point a label from $0, \dots, 7$, and form patch signatures as multisets of labels inside all $16 \times 16$ patches, which gives 38346 patches, of 6956 different unique values. \\
Entropy of resulting categorical distribution of patch signatures: \textbf{12.307} \\

\noindent \begin{large} \textbf{Layer 4} \end{large} \\
Number of points to be clustered: 38346 \\
Number of components found by MDL, as per \eqref{eq:k-star}: 7 \\
Assign each point a label from $0, \dots, 6$, and form patch signatures as multisets of labels inside all size $32 \times 32$ patches, which gives 26666 patches, of 5325 different unique values. \\
Entropy of resulting categorical distribution of patch signatures: \textbf{12.378} \\
\\
\begin{large}
Total complexity: $7.995 + 10.194 + 12.772 + 12.753 = \boldsymbol{43.314}$
\\
\end{large}



\subsection{Cifar Car Image}
\begin{figure*}[h]
    \centering
    \begin{tabular}{@{}c@{}c@{}}
    \includegraphics[width=0.53\textwidth]{imgs/cifar_car.png.png} & 
    \end{tabular}
    \caption{ \small Example of a low-resolution, real-world image from CIFAR.}
    \label{fig:imgs/cifar-car-image}
\end{figure*}
\noindent \begin{large} \textbf{Layer 1} \end{large} \\
Num points to be clustered (pixels): 50176 \\
Num components found by MDL: 7 \\
Assign each pixel a label from $0, \dots, 6$, and form patch signatures as multisets of labels inside all patches of size $4 \times 4$. \\
Num patch signatures: 48400 \\
Num unique patch signatures: 304 \\
Entropy of categorical distribution of patch signatures: \textbf{5.6} \\

\noindent \begin{large} \textbf{Layer 2} \end{large} \\
Num points to be clustered: 48400 \\
Num components found by MDL: 6 \\
Assign each point a label from $0, \dots, 5$, and form patch signatures as multisets of labels inside all patches of size $8 \times 8$. \\
Num patch signatures: 44944 \\
Num unique patch signatures: 2615 \\
Entropy of categorical distribution of patch signatures: \textbf{8.67} \\

\noindent \begin{large} \textbf{Layer 3} \end{large} \\
Num points to be clustered: 44944 \\
Num components found by MDL: 8 \\
Assign each pixel a label from $0, \dots, 7$, and form patch signatures as multisets of labels inside all patches of size $16 \times 16$. \\
Num patch signatures: 38416 \\
Num unique patch signatures: 6455 \\
Entropy of categorical distribution of patch signatures: \textbf{11.69} \\

\noindent \begin{large} \textbf{Layer 4} \end{large} \\
Num points to be clustered: 38212 \\
Num components found by MDL: 8 \\
Assign each pixel a label from $0, \dots, 7$, and form patch signatures as multisets of labels inside all patches of size $32 \times 32$. \\
Num patch signatures: 26896 \\
Num unique patch signatures: 6445 \\
Entropy of categorical distribution of patch signatures: \textbf{12.61} \\
\\
\begin{large}
\\
Total complexity: $5.61 + 8.67 + 11.69 + 12.61 = \boldsymbol{38.50}$
\end{large}
\clearpage

\subsection{DTD Fine Woven Texture}
\begin{figure*}[h]
    \centering
    \begin{tabular}{@{}c@{}c@{}}
    \includegraphics[width=0.53\textwidth]{imgs/woven_dtd.png} & 
    \end{tabular}
    \caption{ \small Example of a detailed, high-resolution, repetitive pattern from DTD2.}
    \label{fig:imgs/dtd-fine-woven-image}
\end{figure*}
\noindent \begin{large} \textbf{Layer 1} \end{large} \\
Num points to be clustered (pixels): 50290 \\
Num components found by MDL: 8 \\
Assign each pixel a label from $0, \dots, 7$, and form patch signatures as multisets of labels inside all patches of size $4 \times 4$. \\
Num patch signatures: 48400 \\
Num unique patch signatures: 7568 \\
Entropy of categorical distribution of patch signatures: \textbf{12.16} \\

\noindent \begin{large} \textbf{Layer 2} \end{large} \\
Num points to be clustered: 48400 \\
Num components found by MDL: 8 \\
Assign each point a label from $0, \dots, 7$, and form patch signatures as multisets of labels inside all patches of size $8 \times 8$. \\
Num patch signatures: 44944 \\
Num unique patch signatures: 13356 \\
Entropy of categorical distribution of patch signatures: \textbf{13.70} \\

\noindent \begin{large} \textbf{Layer 3} \end{large} \\
Num points to be clustered: 44944 \\
Num components found by MDL: 1 \\
Assign each pixel a label from $0, \dots, 7$, and form patch signatures as multisets of labels inside all patches of size $16 \times 16$. \\
Num patch signatures: 36416 \\
Num unique patch signatures: 1 \\
Entropy of categorical distribution of patch signatures: \textbf{0} \\

\noindent \begin{large} \textbf{Layer 4} \end{large} \\
All patch signatures are identical because only one cluster was found at the previous level. So the entropy is \textbf{0}.
\begin{large}
\\
\\
Total complexity: $12.16 + 13.70 + 0 + 0 = \boldsymbol{25.87}$
\end{large}
\clearpage



\section{Datasets: Further Details}
\subsection{Synthetic Datasets}
As described in Section \ref{sec:experimental-evaluation}, I created three synthetic datasets to help test my image complexity metric on a variety of images. The experimental results I report use 500 images sampled from these synthetic datasets. Here, I give the full details for the creation of these datasets. Code will also be released on publication.

\paragraph{Stripes}
These images depict a repeated striped black-and-white pattern. The thickness of the lines, in pixels, is sampled uniformly at random from $[3,10]$, and the slope of the lines is sampled uniformly at random from $[-0.5, -1.5]$. It is sufficient to consider negative slopes only as our method, and all methods that I compare to, are invariant to reflections, so the striped images with slope in $[0.5, 1.5]$ would receive identical scores to those in $[-0.5, -1.5]$. Note that our method is not necessarily invariant to rotations, because it is based on square, axis-aligned patches of pixels. The same is true of the fractal dimension computed with the Minkowski-Bouligand dimension (i.e., the fractal dimension), as it uses a box-counting method. Examples of Stripes images are shown in Figure \ref{fig:imgs/stripes-images}.
\begin{figure}[h]
    \centering
    \begin{tabular}{@{}c@{}c@{}}
    \includegraphics[width=0.53\textwidth]{imgs/stripes_image_curated.png} & 
    \includegraphics[width=0.53\textwidth]{imgs/stripes_image2.png}\\
    \includegraphics[width=0.53\textwidth]{imgs/stripes_image4.png} & 
    \includegraphics[width=0.53\textwidth]{imgs/stripes_image5.png}\\
     \includegraphics[width=0.53\textwidth]{imgs/stripes_image_curated2.png} & 
    \includegraphics[width=0.53\textwidth]{imgs/stripes_image8.png}\\
    \end{tabular}
    \caption{ \small Examples of images from our synthetic Stripes dataset. Most existing methods assign theses images a high complexity. Ours assigns them low, but non-zero complexity.}
    \label{fig:imgs/stripes-images}
\end{figure}



\paragraph{Halves}
These images have one half entirely black and the other entirely white, with the dividing line being at various angles. As with Stripes, the slope of this dividing line is sampled uniformly at random from $[-0.5, -1.5]$. Examples are shown in Figure \ref{fig:halves-images}.
\begin{figure}[h]
    \centering
    \begin{tabular}{@{}c@{}c@{}}
    \includegraphics[width=0.53\textwidth]{imgs/halves_image0.png} & 
    \includegraphics[width=0.53\textwidth]{imgs/halves_image1.png}\\
    \includegraphics[width=0.53\textwidth]{imgs/halves_image2.png} & 
    \includegraphics[width=0.53\textwidth]{imgs/halves_image3.png}\\
     \includegraphics[width=0.53\textwidth]{imgs/halves_image4.png} & 
    \includegraphics[width=0.53\textwidth]{imgs/halves_image5.png}\\
    \end{tabular}
    \caption{ \small Examples of images from our synthetic Halves dataset. Our method assigns these low complexity as do existing methods. However, when we break the method down by scale, as discussed in Section 4.2, we see that it assigns some complexity at a high scale, more so than, e.g., Stripes, because there is some difference between different parts at a high scale, whereas in Stripes, both halves of each image are the same.}
    \label{fig:halves-images}
\end{figure}


\paragraph{Rand}
These images are white noise. Their values are sampled uniformly at random from $[0,1]$, independently for each location and each of three colour channels. Examples are shown in Figure \ref{fig:imgs/rand-images}.

\begin{figure*}[h]
    \centering
    \begin{tabular}{@{}c@{}c@{}}
    \includegraphics[width=0.53\textwidth]{imgs/rand_image0.png} & 
    \includegraphics[width=0.53\textwidth]{imgs/rand_image1.png}\\
    \includegraphics[width=0.53\textwidth]{imgs/rand_image2.png} & 
    \includegraphics[width=0.53\textwidth]{imgs/rand_image3.png}\\
     \includegraphics[width=0.53\textwidth]{imgs/rand_image4.png} & 
    \includegraphics[width=0.53\textwidth]{imgs/rand_image5.png}\\
    \end{tabular}
    \caption{ \small Examples of the white noise images I used in the Rand dataset. Existing complexity measures assign these images a very high complexity. Our method, in contrast, gives them all a zero complexity.}
    \label{fig:imgs/rand-images}
\end{figure*}


\clearpage

\subsection{Existing Datasets}
For the experiments in Section 4, I randomly sample 500 images from each of ImageNet, CiFAR, and MNIST. For DTD2, we manually search through all 5640 images in the original Describable Textures Dataset \cite{cimpoi2014describing}, and we find 341 images with fine detailed but repetitive textures. This section contains further information on the images used for each dataset.

\paragraph{ImageNet}
All images are from the Imagenette subset of ImageNet. The list of labels that we present is taken from the Imagenette labels, available at https://s3.amazonaws. com/fast-ai-imageclas/imagenette2.tgz. Most of the images IDs are of the form <wordnet-synset-id> - <index-within-synset>. Some instead use the class label from the 2012 version of ILSVRC.

\begin{multicols}{3}

\begin{itemize}

\item n03394916 - 50642.
\item n02979186 - 10250.
\item n03028079 - 14492.
\item n03888257 - 36631.
\item n03445777 - 3291.
\item n03425413 - 14940.
\item n03417042 - 19472.
\item n03000684 - 9440.
\item n01440764 - 16090.
\item ILSVRC2012-00023440.
\item n03394916 - 39102.
\item n02979186 - 24592.
\item n03028079 - 25712.
\item n03888257 - 16542.
\item n03445777 - 12262.
\item n03425413 - 17220.
\item n03417042 - 3351.
\item n03000684 - 32351.
\item n01440764 - 21191.
\item n02102040 - 1782.
\item n03394916 - 46700.
\item n02979186 - 17680.
\item n03028079 - 7422.
\item n03888257 - 25150.
\item n03445777 - 5582.
\item n03425413 - 8801.
\item n03417042 - 25411.
\item n03000684 - 20052.
\item n01440764 - 5361.
\item n02102040 - 2930.
\item n03394916 - 43381.
\item n02979186 - 20620.
\item n03028079 - 16811.
\item n03888257 - 21201.
\item n03445777 - 11171.
\item n03425413 - 13581.
\item n03417042 - 6420.
\item ILSVRC2012-00045501.
\item n01440764 - 16192.
\item n02102040 - 5890.
\item n03394916 - 43532.
\item n02979186 - 9910.
\item ILSVRC2012-00004912.
\item n03888257 - 22330.
\item n03445777 - 10762.
\item n03425413 - 11061.
\item n03417042 - 9620.
\item n03000684 - 16872.
\item n01440764 - 6812.
\item n02102040 - 3452.
\item n03394916 - 32870.
\item n02979186 - 23650.
\item n03028079 - 27781.
\item n03888257 - 20300.
\item n03445777 - 6091.
\item n03425413 - 20371.
\item n03417042 - 3771.
\item n03000684 - 10992.
\item n01440764 - 13842.
\item n02102040 - 6851.
\item n03394916 - 11582.
\item n02979186 - 18720.
\item n03028079 - 29062.
\item n03888257 - 18441.
\item n03445777 - 520.
\item n03425413 - 21180.
\item n03417042 - 4072.
\item n03000684 - 9452.
\item n01440764 - 14150.
\item n02102040 - 652.
\item n03394916 - 36000.
\item n02979186 - 1810.
\item n03028079 - 4612.
\item n03888257 - 7921.
\item n03445777 - 5932.
\item n03425413 - 12711.
\item n03417042 - 4462.
\item n03000684 - 661.
\item n01440764 - 9152.
\item n02102040 - 7942.
\item n03394916 - 36172.
\item n02979186 - 13740.
\item n03028079 - 9920.
\item n03888257 - 36390.
\item n03445777 - 6162.
\item n03425413 - 12951.
\item n03417042 - 2150.
\item n03000684 - 10690.
\item n01440764 - 14342.
\item n02102040 - 5942.
\item n03394916 - 43422.
\item n02979186 - 13442.
\item n03028079 - 34051.
\item n03888257 - 46870.
\item n03445777 - 13480.
\item n03425413 - 20751.
\item n03417042 - 5920.
\item n03000684 - 18020.
\item n01440764 - 7982.
\item n02102040 - 182.
\item n03394916 - 42721.
\item n02979186 - 11971.
\item n03028079 - 10191.
\item n03888257 - 19580.
\item n03445777 - 1390.
\item n03425413 - 13862.
\item n03417042 - 18582.
\item ILSVRC2012-00045940.
\item n01440764 - 1561.
\item n02102040 - 4090.
\item n03394916 - 33380.
\item n02979186 - 2002.
\item n03028079 - 9682.
\item n03888257 - 9770.
\item n03445777 - 11162.
\item n03425413 - 15321.
\item n03417042 - 14000.
\item n03000684 - 15441.
\item n01440764 - 16051.
\item n02102040 - 6552.
\item n03394916 - 59361.
\item n02979186 - 9811.
\item n03028079 - 29942.
\item n03888257 - 20352.
\item n03445777 - 2611.
\item n03425413 - 11180.
\item n03417042 - 26782.
\item n03000684 - 7222.
\item n01440764 - 19302.
\item ILSVRC2012-00036282.
\item n03394916 - 62451.
\item n02979186 - 140.
\item n03028079 - 49281.
\item n03888257 - 14530.
\item n03445777 - 5240.
\item n03425413 - 21730.
\item n03417042 - 12790.
\item n03000684 - 13402.
\item n01440764 - 4360.
\item n02102040 - 352.
\item n03394916 - 46672.
\item n02979186 - 1542.
\item n03028079 - 15392.
\item n03888257 - 10680.
\item n03445777 - 17492.
\item n03425413 - 16220.
\item n03417042 - 7080.
\item n03000684 - 16291.
\item n01440764 - 2921.
\item n02102040 - 8061.
\item n03394916 - 30072.
\item n02979186 - 5321.
\item n03028079 - 17690.
\item n03888257 - 70632.
\item n03445777 - 9572.
\item n03425413 - 1672.
\item n03417042 - 4761.
\item n03000684 - 18591.
\item n01440764 - 8030.
\item n02102040 - 5641.
\item n03394916 - 50730.
\item n02979186 - 8861.
\item ILSVRC2012-00016542.
\item n03888257 - 3651.
\item n03445777 - 5312.
\item n03425413 - 21362.
\item n03417042 - 8822.
\item n03000684 - 19272.
\item n01440764 - 6421.
\item n02102040 - 960.
\item n03394916 - 26422.
\item n02979186 - 3260.
\item n03028079 - 6110.
\item n03888257 - 33021.
\item n03445777 - 15810.
\item n03425413 - 8661.
\item n03417042 - 21361.
\item n03000684 - 2820.
\item n01440764 - 650.
\item n02102040 - 1791.
\item n03394916 - 52191.
\item n02979186 - 14630.
\item n03028079 - 6722.
\item n03888257 - 142.
\item n03445777 - 11150.
\item n03425413 - 20500.
\item n03417042 - 27630.
\item n03000684 - 15521.
\item n01440764 - 6130.
\item n02102040 - 491.
\item n03394916 - 71910.
\item n02979186 - 8092.
\item n03028079 - 5942.
\item n03888257 - 11222.
\item n03445777 - 2530.
\item n03425413 - 602.
\item n03417042 - 5221.
\item n03000684 - 1970.
\item n01440764 - 13702.
\item n02102040 - 3450.
\item n03394916 - 35320.
\item n02979186 - 16142.
\item n03028079 - 14992.
\item n03888257 - 37362.
\item n03445777 - 6042.
\item n03425413 - 12712.
\item n03417042 - 26850.
\item n03000684 - 180.
\item n01440764 - 12881.
\item n02102040 - 4111.
\item n03394916 - 16601.
\item n02979186 - 4511.
\item n03028079 - 5432.
\item n03888257 - 64711.
\item n03445777 - 7711.
\item n03425413 - 17212.
\item n03417042 - 5510.
\item n03000684 - 19890.
\item n01440764 - 27422.
\item n02102040 - 651.
\item n03394916 - 54570.
\item n02979186 - 11.
\item n03028079 - 16731.
\item n03888257 - 13410.
\item n03445777 - 7090.
\item n03425413 - 13970.
\item n03417042 - 27862.
\item n03000684 - 2340.
\item n01440764 - 3782.
\item n02102040 - 290.
\item n03394916 - 59430.
\item n02979186 - 26820.
\item n03028079 - 3600.
\item n03888257 - 12401.
\item n03445777 - 1750.
\item n03425413 - 14302.
\item n03417042 - 28552.
\item n03000684 - 11511.
\item n01440764 - 20451.
\item n02102040 - 371.
\item n03394916 - 47852.
\item n02979186 - 3472.
\item n03028079 - 8572.
\item n03888257 - 14901.
\item n03445777 - 3301.
\item n03425413 - 14510.
\item n03417042 - 2141.
\item n03000684 - 31112.
\item n01440764 - 2102.
\item n02102040 - 2572.
\item n03394916 - 38680.
\item n02979186 - 1200.
\item n03028079 - 17922.
\item n03888257 - 15382.
\item n03445777 - 13462.
\item n03425413 - 20121.
\item n03417042 - 15592.
\item n03000684 - 31721.
\item n01440764 - 32420.
\item n02102040 - 1830.
\item n03394916 - 35811.
\item n02979186 - 12072.
\item n03028079 - 46322.
\item n03888257 - 28581.
\item n03445777 - 602.
\item n03425413 - 32871.
\item n03417042 - 18042.
\item n03000684 - 6220.
\item n01440764 - 17501.
\item n02102040 - 7392.
\item n03394916 - 36361.
\item n02979186 - 22761.
\item n03028079 - 24471.
\item n03888257 - 13790.
\item n03445777 - 7930.
\item n03425413 - 21040.
\item n03417042 - 1330.
\item n03000684 - 1542.
\item n01440764 - 8302.
\item n02102040 - 6081.
\item n03394916 - 27071.
\item n02979186 - 5781.
\item ILSVRC2012-00034021.
\item n03888257 - 38102.
\item n03445777 - 16321.
\item n03425413 - 20562.
\item n03417042 - 4560.
\item n03000684 - 6471.
\item n01440764 - 762.
\item n02102040 - 2110.
\item n03394916 - 44882.
\item n02979186 - 5481.
\item n03028079 - 9220.
\item n03888257 - 19211.
\item n03445777 - 14301.
\item n03425413 - 19050.
\item n03417042 - 6691.
\item n03000684 - 2972.
\item n01440764 - 10040.
\item n02102040 - 430.
\item n03394916 - 46391.
\item n02979186 - 13281.
\item n03028079 - 16820.
\item n03888257 - 30712.
\item n03445777 - 14232.
\item n03425413 - 21562.
\item n03417042 - 29412.
\item n03000684 - 13182.
\item n01440764 - 10852.
\item n02102040 - 5101.
\item n03394916 - 29940.
\item n02979186 - 2841.
\item n03028079 - 23280.
\item n03888257 - 23192.
\item n03445777 - 2041.
\item n03425413 - 14570.
\item n03417042 - 20280.
\item n03000684 - 8411.
\item n01440764 - 7492.
\item n02102040 - 6532.
\item n03394916 - 28590.
\item n02979186 - 560.
\item n03028079 - 38692.
\item n03888257 - 23571.
\item n03445777 - 13680.
\item ILSVRC2012-00000732.
\item n03417042 - 18551.
\item n03000684 - 34440.
\item n01440764 - 522.
\item ILSVRC2012-00008162.
\item n03394916 - 1091.
\item n02979186 - 10151.
\item n03028079 - 12802.
\item n03888257 - 171.
\item n03445777 - 7670.
\item n03425413 - 21202.
\item n03417042 - 9601.
\item ILSVRC2012-00029211.
\item n01440764 - 5432.
\item n02102040 - 4732.
\item n03394916 - 292.
\item n02979186 - 5460.
\item n03028079 - 3700.
\item n03888257 - 35800.
\item n03445777 - 9921.
\item n03425413 - 21911.
\item n03417042 - 5090.
\item n03000684 - 19211.
\item n01440764 - 8601.
\item n02102040 - 7841.
\item n03394916 - 27932.
\item n02979186 - 3161.
\item n03028079 - 29012.
\item n03888257 - 17340.
\item n03445777 - 10782.
\item n03425413 - 11161.
\item n03417042 - 29722.
\item n03000684 - 1490.
\item n01440764 - 4962.
\item n02102040 - 7792.
\item n03394916 - 47110.
\item n02979186 - 16952.
\item n03028079 - 28242.
\item n03888257 - 29762.
\item n03445777 - 230.
\item n03425413 - 3021.
\item n03417042 - 10462.
\item n03000684 - 2060.
\item n01440764 - 6301.
\item n02102040 - 2480.
\item n03394916 - 44580.
\item n02979186 - 20362.
\item n03028079 - 3492.
\item n03888257 - 30412.
\item n03445777 - 13831.
\item n03425413 - 20301.
\item n03417042 - 10280.
\item n03000684 - 16861.
\item n01440764 - 9212.
\item n02102040 - 6152.
\item n03394916 - 34332.
\item n02979186 - 14251.
\item n03028079 - 9320.
\item n03888257 - 35890.
\item n03445777 - 5131.
\item n03425413 - 16221.
\item n03417042 - 5381.
\item n03000684 - 3470.
\item n01440764 - 8142.
\item n02102040 - 762.
\item n03394916 - 51071.
\item n02979186 - 20160.
\item n03028079 - 25542.
\item n03888257 - 57010.
\item n03445777 - 261.
\item n03425413 - 7731.
\item n03417042 - 3821.
\item ILSVRC2012-00047060.
\item n01440764 - 12971.
\item n02102040 - 1300.
\item n03394916 - 7292.
\item n02979186 - 23362.
\item n03028079 - 10020.
\item ILSVRC2012-00038942.
\item n03445777 - 11690.
\item n03425413 - 13100.
\item n03417042 - 6811.
\item n03000684 - 20762.
\item n01440764 - 11350.
\item n02102040 - 1822.
\item n03394916 - 33012.
\item n02979186 - 1061.
\item n03028079 - 16660.
\item n03888257 - 38200.
\item n03445777 - 10671.
\item n03425413 - 6772.
\item n03417042 - 1492.
\item n03000684 - 24991.
\item n01440764 - 7462.
\item n02102040 - 362.
\item n03394916 - 26802.
\item n02979186 - 3530.
\item n03028079 - 80.
\item n03888257 - 66102.
\item n03445777 - 8192.
\item ILSVRC2012-00035211.
\item n03417042 - 10300.
\item n03000684 - 16072.
\item n01440764 - 8451.
\item n02102040 - 3260.
\item ILSVRC2012-00025761.
\item n02979186 - 5031.
\item n03028079 - 10241.
\item n03888257 - 12400.
\item n03445777 - 6201.
\item n03425413 - 260.
\item n03417042 - 2062.
\item n03000684 - 27850.
\item n01440764 - 9491.
\item n02102040 - 821.
\item n03394916 - 32340.
\item n02979186 - 1932.
\item n03028079 - 26291.
\item n03888257 - 9552.
\item n03445777 - 101.
\item n03425413 - 1792.
\item n03417042 - 18152.
\item n03000684 - 5041.
\item n01440764 - 4980.
\item n02102040 - 3532.
\item n03394916 - 6742.
\item n02979186 - 22882.
\item n03028079 - 25462.
\item n03888257 - 8381.
\item n03445777 - 5382.
\item n03425413 - 13232.
\item n03417042 - 9170.
\item n03000684 - 17330.
\item n01440764 - 6361.
\item n02102040 - 142.
\item n03394916 - 51161.
\item n02979186 - 15931.
\item n03028079 - 28662.
\item n03888257 - 12070.
\item n03445777 - 10401.
\item n03425413 - 4511.
\item n03417042 - 1601.
\item n03000684 - 10212.
\item n01440764 - 7752.
\item n02102040 - 1110.
\item n03394916 - 38212.
\item n02979186 - 1621.
\item n03028079 - 2060.
\item n03888257 - 7610.
\item n03445777 - 7902.
\item n03425413 - 21211.
\item n03417042 - 3390.
\item n03000684 - 11821.
\item n01440764 - 8221.
\item n02102040 - 350.
\item n03394916 - 37321.
\item n02979186 - 2312.
\item n03028079 - 16501.
\item n03888257 - 11081.
\item n03445777 - 471.
\item n03425413 - 24461.
\item n03417042 - 6272.
\item n03000684 - 6460.
\item n01440764 - 7160.
\item n02102040 - 3112.
\item n03394916 - 33221.
\item n02979186 - 15972.
\item n03028079 - 9112.
\item n03888257 - 7130.
\item n03445777 - 8861.
\item n03425413 - 14552.
\item n03417042 - 2960.
\item n03000684 - 5231.
\item n01440764 - 16072.
\item n02102040 - 672.
\end{itemize}
\end{multicols}



\paragraph{DTD2}
The Describable Textures Dataset contains 47 classes, grouped according to texture: bumpy, dotted, lined, veined etc. The image ids below are of the form <class-id> - <index-within-class>. Unsurprisingly, most of the suitable images, i.e., those with detailed repeating textures, are from classes such as `woven', `grid', or `wrinkled'. The number of images of each class is given below.

\begin{figure}[H]
    \centering
    \includegraphics[width=\textwidth]{imgs/dtd2_class_counts.png}
    \vspace{-16pt}
    \caption{Number of images of each class from DTD that we select to be part of our curated dataset, DTD2. We select images that show a detailed, repeating pattern, meaning some types of textures are much more likely to be selected.}
    \label{fig:dtd2-class-counts}
\end{figure}

\begin{multicols}{2}
\begin{itemize}
\item perforated - 0074.
\item meshed - 0164.
\item dotted - 0164.
\item sprinkled - 0066.
\item porous - 0117.
\item woven - 0106.
\item knitted - 0185.
\item crosshatched - 0081.
\item pleated - 0163.
\item banded - 0046.
\item wrinkled - 0129.
\item banded - 0055.
\item braided - 0008.
\item grooved - 0119.
\item dotted - 0131.
\item wrinkled - 0015.
\item bumpy - 0098.
\item woven - 0004.
\item zigzagged - 0109.
\item matted - 0136.
\item stratified - 0174.
\item grooved - 0051.
\item perforated - 0014.
\item grooved - 0089.
\item woven - 0025.
\item wrinkled - 0106.
\item lined - 0159.
\item banded - 0008.
\item matted - 0070.
\item lined - 0109.
\item dotted - 0192.
\item fibrous - 0138.
\item matted - 0150.
\item pleated - 0142.
\item grid - 0088.
\item blotchy - 0038.
\item chequered - 0043.
\item banded - 0081.
\item wrinkled - 0063.
\item waffled - 0156.
\item grid - 0073.
\item grid - 0016.
\item lacelike - 0078.
\item matted - 0071.
\item chequered - 0050.
\item wrinkled - 0132.
\item porous - 0149.
\item stained - 0119.
\item knitted - 0118.
\item pleated - 0116.
\item stained - 0066.
\item knitted - 0144.
\item chequered - 0062.
\item grooved - 0085.
\item blotchy - 0091.
\item knitted - 0150.
\item grid - 0124.
\item pleated - 0168.
\item zigzagged - 0133.
\item grooved - 0058.
\item zigzagged - 0008.
\item stained - 0132.
\item blotchy - 0088.
\item bumpy - 0067.
\item grid - 0049.
\item woven - 0062.
\item blotchy - 0059.
\item matted - 0069.
\item lined - 0133.
\item woven - 0075.
\item bubbly - 0097.
\item matted - 0073.
\item porous - 0151.
\item blotchy - 0083.
\item chequered - 0052.
\item wrinkled - 0041.
\item lacelike - 0096.
\item matted - 0085.
\item fibrous - 0150.
\item banded - 0122.
\item waffled - 0124.
\item fibrous - 0164.
\item grid - 0083.
\item fibrous - 0193.
\item dotted - 0060.
\item meshed - 0176.
\item woven - 0043.
\item woven - 0088.
\item stained - 0090.
\item wrinkled - 0045.
\item pleated - 0090.
\item zigzagged - 0085.
\item veined - 0135.
\item dotted - 0132.
\item stratified - 0046.
\item woven - 0061.
\item woven - 0028.
\item swirly - 0074.
\item matted - 0065.
\item sprinkled - 0065.
\item waffled - 0068.
\item grooved - 0048.
\item perforated - 0066.
\item grid - 0022.
\item woven - 0053.
\item porous - 0152.
\item fibrous - 0160.
\item woven - 0055.
\item matted - 0148.
\item pitted - 0134.
\item flecked - 0060.
\item lacelike - 0020.
\item grid - 0052.
\item woven - 0067.
\item knitted - 0192.
\item flecked - 0053.
\item chequered - 0054.
\item chequered - 0088.
\item lined - 0027.
\item stained - 0030.
\item knitted - 0146.
\item grid - 0078.
\item blotchy - 0070.
\item swirly - 0060.
\item perforated - 0057.
\item porous - 0098.
\item wrinkled - 0087.
\item blotchy - 0096.
\item grooved - 0081.
\item wrinkled - 0039.
\item lined - 0041.
\item flecked - 0126.
\item lined - 0038.
\item dotted - 0135.
\item pitted - 0036.
\item wrinkled - 0103.
\item wrinkled - 0034.
\item grid - 0050.
\item bubbly - 0083.
\item woven - 0001.
\item knitted - 0116.
\item pleated - 0069.
\item wrinkled - 0043.
\item woven - 0059.
\item knitted - 0079.
\item matted - 0128.
\item lacelike - 0017.
\item fibrous - 0165.
\item wrinkled - 0088.
\item grid - 0129.
\item blotchy - 0090.
\item wrinkled - 0017.
\item sprinkled - 0038.
\item woven - 0032.
\item flecked - 0074.
\item woven - 0029.
\item knitted - 0130.
\item crosshatched - 0092.
\item lacelike - 0065.
\item knitted - 0141.
\item grid - 0011.
\item porous - 0099.
\item woven - 0039.
\item woven - 0113.
\item fibrous - 0211.
\item sprinkled - 0067.
\item wrinkled - 0125.
\item crosshatched - 0093.
\item dotted - 0154.
\item woven - 0130.
\item veined - 0075.
\item meshed - 0181.
\item fibrous - 0103.
\item fibrous - 0183.
\item woven - 0082.
\item woven - 0099.
\item perforated - 0041.
\item grid - 0099.
\item grooved - 0084.
\item meshed - 0162.
\item wrinkled - 0036.
\item banded - 0147.
\item porous - 0157.
\item wrinkled - 0108.
\item dotted - 0185.
\item grid - 0089.
\item grid - 0101.
\item woven - 0048.
\item grid - 0066.
\item bumpy - 0190.
\item matted - 0166.
\item woven - 0104.
\item waffled - 0171.
\item wrinkled - 0040.
\item flecked - 0135.
\item swirly - 0151.
\item stratified - 0115.
\item perforated - 0045.
\item woven - 0026.
\item fibrous - 0111.
\item swirly - 0065.
\item perforated - 0026.
\item banded - 0107.
\item woven - 0068.
\item banded - 0037.
\item fibrous - 0204.
\item wrinkled - 0026.
\item waffled - 0178.
\item woven - 0108.
\item grooved - 0164.
\item woven - 0021.
\item fibrous - 0127.
\item banded - 0141.
\item scaly - 0131.
\item woven - 0123.
\item braided - 0167.
\item woven - 0046.
\item grooved - 0057.
\item perforated - 0024.
\item swirly - 0137.
\item grid - 0081.
\item bubbly - 0118.
\item grooved - 0108.
\item wrinkled - 0079.
\item flecked - 0003.
\item fibrous - 0120.
\item wrinkled - 0114.
\item woven - 0083.
\item fibrous - 0110.
\item wrinkled - 0067.
\item lined - 0169.
\item wrinkled - 0025.
\item wrinkled - 0021.
\item wrinkled - 0013.
\item dotted - 0041.
\item woven - 0049.
\item lined - 0076.
\item scaly - 0122.
\item grid - 0059.
\item waffled - 0081.
\item matted - 0117.
\item fibrous - 0101.
\item stained - 0075.
\item woven - 0036.
\item wrinkled - 0086.
\item wrinkled - 0084.
\item banded - 0047.
\item banded - 0068.
\item matted - 0155.
\item perforated - 0080.
\item pitted - 0078.
\item pitted - 0008.
\item fibrous - 0089.
\item sprinkled - 0068.
\item woven - 0084.
\item grooved - 0093.
\item woven - 0056.
\item pitted - 0064.
\item wrinkled - 0046.
\item woven - 0109.
\item banded - 0086.
\item grid - 0084.
\item grid - 0116.
\item woven - 0038.
\item pleated - 0082.
\item bumpy - 0140.
\item wrinkled - 0111.
\item matted - 0115.
\item fibrous - 0096.
\item woven - 0093.
\item swirly - 0144.
\item banded - 0059.
\item lined - 0141.
\item woven - 0003.
\item banded - 0114.
\item woven - 0002.
\item woven - 0127.
\item knitted - 0126.
\item banded - 0115.
\item woven - 0051.
\item lined - 0166.
\item bubbly - 0084.
\item flecked - 0165.
\item wrinkled - 0105.
\item woven - 0114.
\item woven - 0126.
\item grooved - 0063.
\item wrinkled - 0083.
\item grooved - 0088.
\item grid - 0032.
\item wrinkled - 0065.
\item grid - 0067.
\item perforated - 0016.
\item meshed - 0108.
\item blotchy - 0082.
\item grooved - 0045.
\item swirly - 0147.
\item grooved - 0068.
\item woven - 0065.
\item grid - 0085.
\item blotchy - 0089.
\item wrinkled - 0085.
\item woven - 0107.
\item stratified - 0100.
\item fibrous - 0201.
\item scaly - 0137.
\item woven - 0071.
\item perforated - 0119.
\item swirly - 0138.
\item grooved - 0083.
\item matted - 0072.
\item grid - 0093.
\item chequered - 0093.
\item knitted - 0098.
\item matted - 0084.
\item pitted - 0010.
\item porous - 0053.
\item matted - 0129.
\item woven - 0066.
\item crosshatched - 0109.
\item wrinkled - 0033.
\item pitted - 0157.
\item porous - 0142.
\item woven - 0092.
\item crosshatched - 0116.
\item swirly - 0159.
\item grid - 0082.
\item banded - 0099.
\item meshed - 0112.
\item knitted - 0155.
\item woven - 0112.
\item meshed - 0161.
\item banded - 0061.
\item bubbly - 0055.
\item woven - 0063.
\item bubbly - 0096.
\item blotchy - 0041.
\item banded - 0002.
\item woven - 0007.
\item banded - 0090.
\item pleated - 0094.
\end{itemize}
\end{multicols}
\paragraph{CIFAR10 and MNIST}
For CIFAR and MNIST, we report the distribution of classes in our random sample of 500, in Tables \ref{tab:cifar-counts} and \ref{tab:mnist-counts}, respectively.

\begin{table}
    \caption{Counts of each class in our random sample of 500 CIFAR images.}
    
    \smallskip 
    \label{tab:cifar-counts}
    \centering
    \begin{tabular}{c|*{10}{c}}
    \toprule
        \textbf{class} & airplane & car & bird & cat & deer & dog & frog & horse & ship & truck \\
         \midrule
        \textbf{count} &  66 & 66 & 48 & 48 & 48 & 47 & 49 & 40 & 48 & 40 \\
        \bottomrule
    \end{tabular}
    \end{table}

\begin{table}
    \caption{Counts of each class in our random sample of 500 MNIST images.}
    
    \smallskip 
    \label{tab:mnist-counts}
    \centering
    \begin{tabular}{c|*{10}{c}}
    \toprule
           \textbf{class} & 0 & 1 & 2 & 3 & 4 & 5 & 6 & 7 & 8 & 9 \\
         \midrule
        \textbf{count} &  47 & 61 & 57 & 52 & 47 & 45 & 41 & 50 & 53 & 47 \\
        \bottomrule
    \end{tabular}
    \end{table}
